\chapter{JSON, TOML, YAML, and CSV}\label{ch:data-formats}

\index{JSON}
\index{TOML}
\index{YAML}
\index{CSV}
\index{data formats}

Programs rarely exist in isolation. They read configuration files, consume
API responses, export reports, and exchange data with other programs. The
format of that data matters, and there is no single format that fits every
situation. JSON dominates web APIs. TOML is the configuration language of
choice for many modern tools. YAML appears in CI pipelines and Kubernetes
manifests. CSV is the lingua franca of spreadsheets and data science.

Lattice provides built-in support for all four. Each format has a
\lstinline|_parse| function that turns a string into Lattice values, and a
\lstinline|_stringify| function that turns Lattice values back into a string.
The mapping between formats and Lattice types is consistent and predictable:
JSON objects become Maps, arrays become Arrays, and scalars become the
corresponding Lattice primitives.

Let's start with JSON, the format you will probably use most often.

% ---------------------------------------------------------------------------
\section{json\_parse / json\_stringify}\label{sec:json}
% ---------------------------------------------------------------------------

\index{json\_parse}
\index{json\_stringify}

\subsection{Parsing JSON}

\begin{lstlisting}[language=Lattice, caption={Parsing a JSON string}]
let text = "{\"name\": \"Lattice\", \"version\": 3, \"stable\": true}"
let data = json_parse(text)

print(data["name"])     // "Lattice"
print(data["version"])  // 3
print(data["stable"])   // true
\end{lstlisting}

\lstinline|json_parse| takes a \lstinline|String| and returns a Lattice value.
The type mapping follows the JSON specification exactly:

\begin{center}
\begin{tabular}{ll}
\toprule
\textbf{JSON Type} & \textbf{Lattice Type} \\
\midrule
object & \lstinline|Map| \\
array & \lstinline|Array| \\
string & \lstinline|String| \\
number (integer) & \lstinline|Int| \\
number (decimal/exponent) & \lstinline|Float| \\
\lstinline|true| / \lstinline|false| & \lstinline|Bool| \\
\lstinline|null| & \lstinline|nil| \\
\bottomrule
\end{tabular}
\end{center}

Notice that the parser distinguishes between integers and floats: \texttt{42}
becomes an \lstinline|Int|, while \texttt{42.0} or \texttt{4.2e1} becomes a
\lstinline|Float|. This matters because integer arithmetic in Lattice is
exact, while floating-point arithmetic is subject to the usual IEEE 754
rounding.

\begin{lstlisting}[language=Lattice, caption={Parsing nested JSON structures}]
let json_text = """
{
    "users": [
        {"id": 1, "name": "Alice", "active": true},
        {"id": 2, "name": "Bob", "active": false}
    ],
    "count": 2
}
"""

let data = json_parse(json_text)
let users = data["users"]

for user in users {
    let status = if user["active"] { "active" } else { "inactive" }
    print(user["name"] + " is " + status)
}
// Alice is active
// Bob is inactive
\end{lstlisting}

\begin{warnbox}{Error Handling in json\_parse}
If the input is not valid JSON, \lstinline|json_parse| sets a runtime error
with a message indicating the position of the problem. The parser in
\filename{src/json.c} uses a recursive descent approach and validates
strictly---trailing content after the root value is rejected, and all escape
sequences must be well-formed. Wrap parsing calls in \lstinline|try|/\lstinline|catch|
if you are dealing with untrusted input.
\end{warnbox}

\subsection{Stringifying to JSON}

\begin{lstlisting}[language=Lattice, caption={Converting Lattice values to JSON}]
let user = Map::new()
user["name"] = "Charlie"
user["age"] = 28
user["hobbies"] = ["climbing", "cooking", "Lattice"]

let json_text = json_stringify(user)
print(json_text)
// {"name":"Charlie","age":28,"hobbies":["climbing","cooking","Lattice"]}
\end{lstlisting}

\lstinline|json_stringify| accepts any Lattice value that has a natural JSON
representation: Maps, Arrays, Strings, Ints, Floats, Bools, and
\lstinline|nil|. Tuples and Buffers are serialized as JSON arrays. Refs are
unwrapped automatically.

\begin{notebox}{Unsupported Types}
Structs, closures, channels, ranges, sets, iterators, and enums cannot be
serialized to JSON directly. Attempting to stringify one of these will produce
a runtime error. If you need to serialize a struct, convert it to a Map first.
\end{notebox}

The serializer handles special floating-point values gracefully:
\lstinline|Infinity|, \lstinline|-Infinity|, and \lstinline|NaN| are all
encoded as \texttt{null}, since JSON has no representation for them. Strings
are properly escaped, including Unicode sequences (\texttt{\textbackslash
uXXXX}) for control characters below \texttt{U+0020}.

\subsection{A Practical Example: Reading and Writing JSON Files}

\begin{lstlisting}[language=Lattice, caption={Loading and saving application state as JSON}]
fn load_state(path: String) -> Map {
    if !file_exists(path) {
        let defaults = Map::new()
        defaults["theme"] = "dark"
        defaults["font_size"] = 14
        defaults["recent_files"] = []
        return defaults
    }
    let text = read_file(path)
    return json_parse(text)
}

fn save_state(path: String, state: Map) {
    let text = json_stringify(state)
    write_file(path, text)
}

// Load, modify, save
flux state = load_state("settings.json")
state["font_size"] = 16
state["recent_files"] = ["main.lat", "lib/utils.lat"]
save_state("settings.json", state)
\end{lstlisting}

% ---------------------------------------------------------------------------
\section{toml\_parse / toml\_stringify}\label{sec:toml}
% ---------------------------------------------------------------------------

\index{toml\_parse}
\index{toml\_stringify}

TOML (Tom's Obvious Minimal Language) was designed to be a configuration file
format that is easy to read and write by hand. If JSON is for machines, TOML
is for humans.

\subsection{Parsing TOML}

\begin{lstlisting}[language=Lattice, caption={Parsing a TOML configuration}]
let toml_text = """
title = "My Project"
version = 3

[database]
host = "localhost"
port = 5432
enabled = true

[logging]
level = "info"
file = "/var/log/app.log"
"""

let config = toml_parse(toml_text)
print(config["title"])              // "My Project"
print(config["database"]["host"])   // "localhost"
print(config["database"]["port"])   // 5432
print(config["logging"]["level"])   // "info"
\end{lstlisting}

Like \lstinline|json_parse|, the TOML parser returns a Map. Table headers
(\texttt{[section]}) create nested Maps. The parser supports all standard
TOML features:

\begin{itemize}
    \item Bare keys (\texttt{name = "value"}) and quoted keys
          (\texttt{"complex key" = "value"})
    \item Dotted keys (\texttt{server.host = "localhost"}) that create
          nested maps automatically
    \item Single-quoted strings (literal, no escape processing) and
          double-quoted strings (with escape sequences)
    \item Integers (including underscored: \texttt{1\_000\_000}),
          floats, and booleans
    \item Inline arrays (\texttt{[1, 2, 3]}) and inline tables
          (\texttt{\{key = "value"\}})
    \item Array-of-tables syntax (\texttt{[[items]]})
\end{itemize}

\begin{lstlisting}[language=Lattice, caption={TOML array of tables}]
let manifest = """
[package]
name = "my-app"
version = "1.0.0"

[[dependencies]]
name = "http"
version = "0.5.0"

[[dependencies]]
name = "json"
version = "1.2.0"
"""

let pkg = toml_parse(manifest)
print(pkg["package"]["name"])  // "my-app"

let deps = pkg["dependencies"]
for dep in deps {
    print(dep["name"] + " v" + dep["version"])
}
// http v0.5.0
// json v1.2.0
\end{lstlisting}

Each \texttt{[[dependencies]]} block pushes a new Map onto an array under the
\texttt{"dependencies"} key. This is how TOML represents arrays of structured
data.

\subsection{Stringifying to TOML}

\begin{lstlisting}[language=Lattice, caption={Generating TOML from Lattice values}]
let config = Map::new()
config["title"] = "Generated Config"
config["debug"] = false

let server = Map::new()
server["host"] = "0.0.0.0"
server["port"] = 8080
config["server"] = server

let output = toml_stringify(config)
print(output)
// title = "Generated Config"
// debug = false
//
// [server]
// host = "0.0.0.0"
// port = 8080
\end{lstlisting}

The TOML serializer in \filename{src/toml\_ops.c} uses a two-pass approach:
first it emits all scalar key-value pairs for the current table, then it
recurses into sub-tables and arrays of tables. This produces clean,
well-structured TOML output.

\begin{warnbox}{toml\_stringify Requires a Map}
Unlike \lstinline|json_stringify|, which accepts any serializable value,
\lstinline|toml_stringify| requires its argument to be a Map. TOML documents
are always rooted in a table (a Map), so passing an Array or a String will
produce a runtime error.
\end{warnbox}

% ---------------------------------------------------------------------------
\section{yaml\_parse / yaml\_stringify}\label{sec:yaml}
% ---------------------------------------------------------------------------

\index{yaml\_parse}
\index{yaml\_stringify}

YAML is the most flexible of the four formats. It uses indentation for
structure (like Python), supports both block and flow styles, and has
extensive scalar type detection.

\subsection{Parsing YAML}

\begin{lstlisting}[language=Lattice, caption={Parsing a YAML document}]
let yaml_text = """
name: Lattice
version: 3
features:
  - phases
  - structured concurrency
  - pattern matching
database:
  host: localhost
  port: 5432
  ssl: true
"""

let data = yaml_parse(yaml_text)
print(data["name"])               // "Lattice"
print(data["database"]["host"])   // "localhost"
print(data["database"]["ssl"])    // true

for feature in data["features"] {
    print("- " + feature)
}
// - phases
// - structured concurrency
// - pattern matching
\end{lstlisting}

Lattice's YAML parser handles the most commonly used subset of YAML:

\begin{itemize}
    \item Block mappings (\texttt{key: value})
    \item Block sequences (\texttt{- item})
    \item Flow sequences (\texttt{[a, b, c]}) and flow mappings
          (\texttt{\{a: 1, b: 2\}})
    \item Quoted and unquoted strings
    \item Document separators (\texttt{---})
\end{itemize}

\begin{notebox}{Automatic Scalar Detection}
The YAML parser automatically detects scalar types from unquoted values.
The words \texttt{true}, \texttt{True}, \texttt{yes}, and \texttt{Yes} all
become \lstinline|Bool| \lstinline|true|. Similarly, \texttt{false},
\texttt{no}, and their capitalized variants become \lstinline|false|. The
words \texttt{null}, \texttt{Null}, and \texttt{\textasciitilde} become
\lstinline|nil|. Numeric-looking strings are parsed as \lstinline|Int| or
\lstinline|Float|. If none of these patterns match, the value remains a
\lstinline|String|. You can see this logic in the
\texttt{yaml\_detect\_scalar} function in \filename{src/yaml\_ops.c}.
\end{notebox}

\begin{lstlisting}[language=Lattice, caption={YAML scalar auto-detection}]
let yaml_text = """
string_val: hello
int_val: 42
float_val: 3.14
bool_val: true
null_val: null
yes_val: yes
"""

let data = yaml_parse(yaml_text)
print(typeof(data["string_val"]))  // "String"
print(typeof(data["int_val"]))     // "Int"
print(typeof(data["float_val"]))   // "Float"
print(typeof(data["bool_val"]))    // "Bool"
print(typeof(data["null_val"]))    // "Nil"
print(typeof(data["yes_val"]))     // "Bool"
\end{lstlisting}

\begin{tipbox}{Force Strings with Quotes}
If you need a value like \texttt{yes} or \texttt{3.14} to remain a string
rather than being auto-detected as a boolean or float, wrap it in quotes in
your YAML source: \texttt{value: "yes"} or \texttt{pi: "3.14"}.
\end{tipbox}

\subsection{Stringifying to YAML}

\begin{lstlisting}[language=Lattice, caption={Generating YAML from Lattice values}]
let deployment = Map::new()
deployment["apiVersion"] = "apps/v1"
deployment["kind"] = "Deployment"

let metadata = Map::new()
metadata["name"] = "my-app"
metadata["labels"] = Map::new()
metadata["labels"]["app"] = "my-app"
deployment["metadata"] = metadata

let yaml_out = yaml_stringify(deployment)
print(yaml_out)
// apiVersion: apps/v1
// kind: Deployment
// metadata:
//   name: my-app
//   labels:
//     app: my-app
\end{lstlisting}

The YAML serializer produces clean, indented block-style output. It
automatically quotes strings that could be misinterpreted as YAML special
values---so a string containing \texttt{"true"} will be emitted as
\texttt{"true"} rather than bare \texttt{true}. Empty maps and arrays use
the compact flow forms \texttt{\{\}} and \texttt{[]}.

\lstinline|yaml_stringify| accepts Maps and Arrays. Other types produce a
runtime error.

% ---------------------------------------------------------------------------
\section{csv\_parse / csv\_stringify}\label{sec:csv}
% ---------------------------------------------------------------------------

\index{csv\_parse}
\index{csv\_stringify}

CSV is the oldest and most universal data exchange format. It has no
formal specification (though RFC 4180 comes close), but the basic idea is
universal: rows separated by newlines, fields separated by commas.

\subsection{Parsing CSV}

\begin{lstlisting}[language=Lattice, caption={Parsing CSV data}]
let csv_text = "name,age,city\nAlice,30,Portland\nBob,25,Seattle\n"
let rows = csv_parse(csv_text)

for row in rows {
    print(row)
}
// ["name", "age", "city"]
// ["Alice", "30", "Portland"]
// ["Bob", "25", "Seattle"]
\end{lstlisting}

\lstinline|csv_parse| returns an Array of Arrays---each inner array
represents one row, and each element is a \lstinline|String|. Unlike JSON and
YAML, CSV has no type system, so all values come back as strings. If you
need numbers, use \lstinline|parse_int| or \lstinline|parse_float| on the
individual fields.

The parser handles the important edge cases:

\begin{itemize}
    \item \textbf{Quoted fields}: Fields enclosed in double quotes can contain
          commas, newlines, and other special characters.
    \item \textbf{Escaped quotes}: A double quote inside a quoted field is
          escaped by doubling it: \texttt{"He said ""hello"""} produces
          \texttt{He said "hello"}.
    \item \textbf{Windows line endings}: Both \texttt{\textbackslash n} and
          \texttt{\textbackslash r\textbackslash n} are handled.
\end{itemize}

\begin{lstlisting}[language=Lattice, caption={Parsing CSV with quoted fields}]
let csv_text = """
name,bio,score
Alice,"Enjoys hiking, reading",95
Bob,"Said ""hello"" to everyone",87
"""

let rows = csv_parse(csv_text)
print(rows[1][1])  // "Enjoys hiking, reading"
print(rows[2][1])  // "Said \"hello\" to everyone"
\end{lstlisting}

\subsection{Working with Headers}

Since CSV does not distinguish between header rows and data rows, you will
often want to convert the raw arrays into a more convenient structure:

\begin{lstlisting}[language=Lattice, caption={Converting CSV rows to maps}]
fn csv_to_maps(rows: Array) -> Array {
    let headers = rows[0]
    flux results = []

    flux i = 1
    while i < rows.len() {
        let row = rows[i]
        let record = Map::new()
        flux j = 0
        while j < headers.len() {
            if j < row.len() {
                record[headers[j]] = row[j]
            }
            j += 1
        }
        results.push(record)
        i += 1
    }
    return results
}

let csv_text = "name,age,city\nAlice,30,Portland\nBob,25,Seattle\n"
let records = csv_to_maps(csv_parse(csv_text))

for record in records {
    print(record["name"] + " lives in " + record["city"])
}
// Alice lives in Portland
// Bob lives in Seattle
\end{lstlisting}

\subsection{Stringifying to CSV}

\begin{lstlisting}[language=Lattice, caption={Generating CSV output}]
let data = [
    ["product", "price", "quantity"],
    ["Widget", "9.99", "100"],
    ["Gadget", "24.99", "50"],
    ["Doohickey", "4.99", "200"]
]

let csv_text = csv_stringify(data)
print(csv_text)
// product,price,quantity
// Widget,9.99,100
// Gadget,24.99,50
// Doohickey,4.99,200
\end{lstlisting}

\lstinline|csv_stringify| takes an Array of Arrays and produces a CSV string.
Each inner array becomes a row. Fields that contain commas, double quotes, or
newlines are automatically wrapped in double quotes, with internal quotes
escaped by doubling.

\begin{lstlisting}[language=Lattice, caption={CSV auto-quoting of special characters}]
let data = [
    ["name", "description"],
    ["Widget", "A small, useful device"],
    ["Gadget", "Says \"beep\" loudly"]
]

let csv_text = csv_stringify(data)
print(csv_text)
// name,description
// Widget,"A small, useful device"
// Gadget,"Says ""beep"" loudly"
\end{lstlisting}

% ---------------------------------------------------------------------------
\section{Round-Tripping Data Between Formats}\label{sec:round-tripping}
% ---------------------------------------------------------------------------

\index{round-tripping}

One of the advantages of Lattice's consistent parsing model is that data
flows naturally between formats. Parse from one format, stringify to another.
The intermediate representation---Maps, Arrays, and primitives---is the same
regardless of which format you started with.

\begin{lstlisting}[language=Lattice, caption={Converting TOML configuration to JSON}]
let toml_text = read_file("config.toml")
let config = toml_parse(toml_text)

// Convert to JSON for an API or another tool
let json_text = json_stringify(config)
write_file("config.json", json_text)
print("Converted TOML to JSON")
\end{lstlisting}

\begin{lstlisting}[language=Lattice, caption={Converting YAML to TOML}]
let yaml_text = """
server:
  host: 0.0.0.0
  port: 8080
  workers: 4
database:
  url: postgres://localhost/mydb
  pool_size: 10
"""

let config = yaml_parse(yaml_text)
let toml_text = toml_stringify(config)
write_file("config.toml", toml_text)
print(toml_text)
// [server]
// host = "0.0.0.0"
// port = 8080
// workers = 4
//
// [database]
// url = "postgres://localhost/mydb"
// pool_size = 10
\end{lstlisting}

\subsection{CSV to JSON}

CSV-to-JSON conversion is particularly common when preparing data for web
APIs:

\begin{lstlisting}[language=Lattice, caption={Converting CSV to JSON}]
let csv_text = read_file("employees.csv")
let rows = csv_parse(csv_text)

// First row is headers
let headers = rows[0]
flux records = []

flux i = 1
while i < rows.len() {
    let row = rows[i]
    let record = Map::new()
    flux j = 0
    while j < headers.len() {
        let field = if j < row.len() { row[j] } else { "" }
        // Try to parse numeric fields
        if headers[j] == "age" || headers[j] == "salary" {
            record[headers[j]] = parse_int(field)
        } else {
            record[headers[j]] = field
        }
        j += 1
    }
    records.push(record)
    i += 1
}

let json_text = json_stringify(records)
write_file("employees.json", json_text)
\end{lstlisting}

\begin{notebox}{Lossy Round-Trips}
Not all round-trips are lossless. TOML has no \lstinline|null| type, so a
JSON \texttt{null} will be lost or turned into an empty string when
converted to TOML. YAML's auto-detection means that the string
\texttt{"true"} in JSON might become a boolean if you round-trip through
YAML without quoting. CSV loses all type information entirely---everything
becomes strings. Be aware of these edge cases when converting between
formats.
\end{notebox}

% ---------------------------------------------------------------------------
\section{Building Configuration-Driven Programs}\label{sec:config-driven}
% ---------------------------------------------------------------------------

\index{configuration}

Let's apply what we have learned to build a realistic, configuration-driven
application. We will build a log file analyzer that reads its configuration
from a TOML file, processes log data from a CSV, and outputs results in
JSON.

\subsection{The Configuration File}

\begin{lstlisting}[language=Lattice, caption={analyzer.toml --- configuration file}]
// Suppose analyzer.toml contains:
// [input]
// path = "logs/access.csv"
// delimiter = ","
//
// [filters]
// min_status = 400
// methods = ["GET", "POST"]
//
// [output]
// format = "json"
// path = "reports/errors.json"
// pretty = true
\end{lstlisting}

\subsection{The Analyzer}

\begin{lstlisting}[language=Lattice, caption={A configuration-driven log analyzer}]
fn load_config(path: String) -> Map {
    let text = read_file(path)
    return toml_parse(text)
}

fn analyze_logs(config: Map) {
    // Read input
    let input_path = config["input"]["path"]
    let csv_text = read_file(input_path)
    let rows = csv_parse(csv_text)

    if rows.len() < 2 {
        print("No data to analyze")
        return
    }

    let headers = rows[0]
    let min_status = config["filters"]["min_status"]
    let allowed_methods = config["filters"]["methods"]

    // Find column indices
    flux status_col = -1
    flux method_col = -1
    flux path_col = -1
    flux i = 0
    while i < headers.len() {
        if headers[i] == "status" { status_col = i }
        if headers[i] == "method" { method_col = i }
        if headers[i] == "path" { path_col = i }
        i += 1
    }

    // Filter and collect matching rows
    flux errors = []
    i = 1
    while i < rows.len() {
        let row = rows[i]
        let status = parse_int(row[status_col])
        let method = row[method_col]

        if status >= min_status {
            flux method_match = false
            for allowed in allowed_methods {
                if method == allowed {
                    method_match = true
                    break
                }
            }
            if method_match {
                let record = Map::new()
                record["status"] = status
                record["method"] = method
                record["path"] = row[path_col]
                errors.push(record)
            }
        }
        i += 1
    }

    // Build report
    let report = Map::new()
    report["total_errors"] = errors.len()
    report["errors"] = errors

    // Write output
    let output_path = config["output"]["path"]
    let json_text = json_stringify(report)
    write_file(output_path, json_text)

    print("Found " + to_string(errors.len()) + " errors")
    print("Report written to " + output_path)
}

let config = load_config("analyzer.toml")
analyze_logs(config)
\end{lstlisting}

This program demonstrates the power of combining formats: TOML for
human-readable configuration, CSV for tabular data ingestion, and JSON for
structured output.

\subsection{Using the dotenv Library}

\index{dotenv}

For programs that need environment-variable configuration (common in
twelve-factor apps and containerized deployments), Lattice ships with a
\texttt{dotenv} library that parses \texttt{.env} files:

\begin{lstlisting}[language=Lattice, caption={Loading environment variables from a .env file}]
import "lib/dotenv" as dotenv

// Load .env from the current directory
dotenv.load()

// Access environment variables
let db_url = env("DATABASE_URL")
let secret = env("SECRET_KEY")
print("Connecting to: " + db_url)
\end{lstlisting}

The dotenv library, found in \filename{lib/dotenv.lat}, supports
double-quoted values with escape sequences, single-quoted literal values,
inline comments, the \texttt{export} prefix, multiline values, and even
variable expansion (\texttt{\$\{VAR\}} syntax). It is a good example of a
real-world configuration parser built entirely in Lattice.

\begin{tipbox}{Layered Configuration}
A robust configuration strategy combines multiple sources with clear
precedence: environment variables override \texttt{.env} files, which
override \texttt{config.toml} defaults. The dotenv library's
\lstinline|load_opts| function respects this pattern---by default it does
\emph{not} override variables that are already set in the environment.
\end{tipbox}

% ---------------------------------------------------------------------------
\section*{Exercises}\label{sec:data-formats-exercises}
% ---------------------------------------------------------------------------

\begin{enumerate}
    \item \textbf{JSON Pretty Printer.} Write a function
    \lstinline|json_pretty(value: any) -> String| that produces indented,
    human-readable JSON output. Use recursion to handle nested maps and
    arrays, adding 2 spaces of indentation per level.

    \item \textbf{Config Merger.} Write a function that takes two TOML
    configuration strings, parses both, and deep-merges them---the second
    configuration's values override the first. Nested maps should be merged
    recursively rather than replaced wholesale.

    \item \textbf{CSV Column Filter.} Write a program that reads a CSV file,
    keeps only the columns specified in a TOML configuration file, and writes
    the filtered CSV to a new file. The config should look like:
    \texttt{columns = ["name", "email", "department"]}.

    \item \textbf{Schema Validator.} Write a function that takes a parsed
    JSON value and a ``schema'' (a Map describing expected types for each key)
    and returns an array of validation errors. For example, the schema
    \lstinline|\{"name": "String", "age": "Int"\}| should report an error if
    \texttt{name} is missing or \texttt{age} is a string.

    \item \textbf{Format Converter CLI.} Write a program that reads from
    standard input in one format and writes to standard output in another.
    Accept the input and output formats as command-line arguments:
    \texttt{lattice convert.lat --from yaml --to json < input.yaml}.
\end{enumerate}

% ---------------------------------------------------------------------------
\section*{What's Next}\label{sec:data-formats-next}
% ---------------------------------------------------------------------------

We can now read files and parse structured data. The next frontier is
communicating over the network. In the next chapter, we will explore
Lattice's HTTP client for making web requests, TCP and TLS networking for
lower-level communication, and how to build a working HTTP server---combining
everything we have learned about strings, data formats, and concurrency.
